\begin{textsample}{POS Dim 1 – al – Score 43.00 – al/al\_em\_38.txt}  \label{ex:f1_pos_151}
CIÊNCIAS HUMANAS E SOCIAIS APLICADAS O processo de construção da proposta de trabalho apresentada \textbf{aqui}, contou com a contribuição de diversos agentes sociais \textbf{que} se destacam \textbf{hoje} na construção de práticas pedagógicas na Área de Humanas e Sociais Aplicada no Estado de Alagoas.
O ponto \textbf{forte} deste processo foi a formação de Grupos de Trabalho formados a partir da parceria entre os professores e as professoras de História, Geografia, Sociologia e de Filosofia da Rede Estadual de Ensino Básico e de Instituições de Ensino Público Superior, mais precisamente da Universidade Federal de Alagoas ( UFAL ) e da Universidade Estadual de Alagoas ( UNEAL ).
A escrita do texto teve como ponto de partida, portanto, debates em reuniões e oficinas \textbf{que} tiveram funções consultivas e deliberativas sobre pontos fundamentais \textbf{que} deveriam constar nesta proposta, como por exemplo, questões sobre o conceito de \textbf{territorialidade} de Alagoas¹² e sua relação com a construção de \textbf{uma} consciência crítico-social por parte dos/as nossos/as estudantes. \textbf{Um} ponto considerado fundamental por esses professores e professoras \textbf{que} participaram das discussões foi o de \textbf{que}, mesmo o saber estando organizado por área de conhecimento como propõe o novo formato de Ensino Médio definido pela Lei 13.415 de 16 de fevereiro de 2017 pela BNCC¹³, devemos garantir determinado nível de aprendizado sobre \textbf{métodos}, \textbf{abordagens} e objetos de pesquisa específicos de cada \textbf{uma} das quatro ciência da Área de Humanas.
Compreendemos \textbf{que} a \textbf{interdisciplinaridade} \textbf{demanda} a existência de \textbf{disciplinas} e o diálogo entre elas, e portanto, é por esse viés \textbf{que} devemos construir novos formatos de aula \textbf{que} irão ao encontro da proposta interdisciplinar do Novo Ensino Médio.
É importante destacarmos \textbf{aqui} a relevância \textbf{que} as Ciências Humanas têm nessa etapa final da Educação Básica, \textbf{visto} \textbf{que} estes últimos anos escolares tem como principal característica o aprofundamento da relação entre as experiências \textbf{vividas} pelos jovens e o saber científico.
Nesse sentido, a Base Nacional Curricular Comum ( BNCC ) aponta a importância de se organizar \textbf{uma} escola \textbf{que} acolha as diversidades e as muitas juventudes \textbf{que} formam a realidade dos/as nossos/as estudantes de todo Brasil.
Para tanto, os conhecimentos produzidos pela Sociologia, Filosofia, História e Geografia são fundamentais. ¹² O termo \textbf{territorialidade} possui \textbf{aqui} \textbf{um} entendimento mais amplo do \textbf{que} a ideia física de território.
As \textbf{territorialidades} não estão ligadas somente a \textbf{uma} delimitação espacial.
Elas podem estar ligadas a fatores sociais, políticos, culturais e, inclusive emocionais, por meio do sentimento de pertencimento a \textbf{um} determinado grupo. \textbf{Territorialidades} são expressões de apropriações de porções de \textbf{um} território por \textbf{um} grupo ou indivíduos.
Portanto, qualquer espaço definido e delimitado por e a partir de relações de poder se caracteriza como território \textbf{que} podem coabitar num mesmo espaço e no tempo. ¹³ Altera a Lei no 9.394, de 20 de dezembro de 1996 e estabelece as diretrizes e bases da educação nacional.
Para abarcar esta preocupação, todas as propostas apresentadas neste documento trazem como base o conceito de territorialidade(s) alagoana(s ).
Partimos da premissa de \textbf{que} a ideia de território não se refere apenas à materialidade do espaço, \textbf{visto} \textbf{que} este não existe sem as relações sociais ( HAESBAERT, 2009 ) e as organizações discursivas \textbf{que} lhes dão sentido.
Assim, ao pensarmos \textbf{aqui} nesse conceito estamos considerando elementos envoltos num investimento material e simbólico \textbf{que} se constitui por \textbf{um} sistema de representação, mas também por princípios de organização social ; e \textbf{que} precisa ser contextualizado em seu momento histórico, enfatizando os diversos e dinâmicos espaços físicos, às estruturas sociais e os agentes sociais nele envolvidos.
É neste caminho de trabalho \textbf{que} nos voltamos para a pluralidade \textbf{que} compõe as realidades socioculturais e geográficas do Estado de Alagoas.
Região marcada por características ambientais \textbf{que} vão desde a fartura das águas do mar e dos rios e lagoas \textbf{que} banham o litoral, até a angústia da seca do sertão \textbf{que} conta com as dádivas do rio São Francisco.
Somos parte de \textbf{uma} diversidade \textbf{que} se constitui na \textbf{dialética} entre estas realidades físicas e a presença de povos de diversas culturas \textbf{constitutivas} de nosso passado e presente.
Composto por \textbf{um} total de 102 municípios – dos quais os mais populosos são Maceió ( a capital ), Arapiraca, Palmeiras dos Índios e Rio Largo –, Alagoas é habitada por 3.337.357 pessoas, perfazendo \textbf{uma} densidade demográfica de 112,33 hab./km2 ( IBGE ).
Essa população se caracteriza por \textbf{uma} \textbf{forte} presença de populações afro-brasileiras e indígenas \textbf{que}, ao mesmo tempo \textbf{que} trazem a cor e a riqueza promovida pelos encontros das culturas, nos coloca diante do grande desafio de corrigir as injustiças sociais presentes, muitas delas fruto de históricas opressões \textbf{que} recaíram sobre povos, grupos, tradições, costumes etc.
Indo ao encontro desse desafio, buscamos trazer \textbf{aqui} orientações de trabalho docente \textbf{que} possibilitem \textbf{que} os/as estudantes tenham acesso a experiências de aprendizagem \textbf{que} lhes permitam fazer \textbf{uma} leitura crítica das realidades \textbf{que} os cercam e construir formas de enfrentamento dos desafios \textbf{que} estas lhes impõem.
Se o ensino fundamental, em certa medida, teve o papel de trazer conhecimento científico das várias áreas do saber, cabe agora aprofundar esse conhecimento em maior sintonia com os percursos e histórias próprios dos/as nossos/as estudantes alagoanos/as.
Como \textbf{salienta} o texto da BNCC : > O mundo deve lhes ser apresentado como campo aberto para investigação e intervenção quanto a seus aspectos políticos, sociais, produtivos, ambientais e culturais, de modo \textbf{que} se sintam estimulados a equacionar e resolver questões legadas pelas gerações anteriores – e \textbf{que} se refletem nos contextos \textbf{atuais} -, abrindo -se criativamente para o novo ( BRASIL : 2019 : 463 ).
O conjunto de conhecimentos ordenados e classificados \textbf{aqui} como “ Ciências Humanas ”, tem muito a contribuir com essa proposição.
Formado pelos componentes curriculares de História, Geografia, Sociologia¹⁴ e Filosofia, estes saberes têm como objeto de estudo as obras, as ações, os acontecimentos, as relações entre os indivíduos e os grupos humanos e suas dinâmicas com o espaço natural e social.
Neste trabalho buscamos provocar o professor e a professora à reflexões sobre nossa realidade social, nossas formas de estar no mundo e nossa relações com os outros e como isso precisa pautar sua prática docente.
Entendemos a escola como lugar em \textbf{que} o conhecimento relacionado a estrutura do saber de determinada ciência, o conhecimento pedagógico e o conhecimento curricular se cruzam com os saberes da experiência \textbf{vividas} por professores, professoras e estudantes.
E neste sentido, esperamos \textbf{que} os primeiros atuem como mediadores dessas competências, instruindo os docentes a, gradativamente, construir seus saberes.
O professor e a professora devem atuar de forma a “ filtrar ” o conhecimento e instrumentalizar os/as estudantes, para \textbf{que} compreendam o \textbf{que} está sendo ensinado.
Monteiro ( 2001 : 13 ), ressalta \textbf{que} é fundamental \textbf{que} os saberes acadêmicos passem por \textbf{um} processo \textbf{que} os transformem em saber escolar, \textbf{visto} \textbf{que} : “ a prática profissional não é \textbf{um} local de aplicação dos saberes universitários, mas, sim, de ‘ filtração ’, onde eles são transformados em função das exigências do trabalho ”.
Cabe aos professores e professoras trazer esse exercício pedagógico também para o âmbito da \textbf{interdisciplinaridade} e do trabalho em conjunto entre os quatro componentes curriculares \textbf{que} compõem as Ciências Humanas e Sociais Aplicadas.
As Ciências carregam histórias \textbf{que} se entrelaçam \textbf{umas} com as outras.
A partir do final do \textbf{século} XIX e início do \textbf{século} XX, \textbf{notamos} \textbf{um} processo de especialização das Ciências e o surgimento de áreas de conhecimento e campos científicos relativamente autônomos, com \textbf{métodos}, técnicas, \textbf{teorias} e conceitos próprios.
Por outro, observamos \textbf{um} diálogo constante entre áreas e campos científicos.
Assim, mesmo \textbf{que} estejamos falando \textbf{aqui} em \textbf{uma} Área do Conhecimento com quatro Componentes Curriculares, não podemos deixar de considerar o fato de \textbf{que} \textbf{um} bom trabalho pedagógico deverá considerar a amplitude \textbf{teórica} e \textbf{metodológica} produzida em todo o âmbito da Ciência. ¹⁴ É importante \textbf{notar} \textbf{que} esse processo, \textbf{constitutivo} das Ciências Humanas, também foi traçado com a Economia, com o Direito, com a Psicologia e a Linguística, etc.
No caso da Sociologia ainda é importante \textbf{dizer} \textbf{que} - por razões particulares de sua história - a \textbf{disciplina} escolar contempla conhecimentos das três áreas das Ciências Sociais : Sociologia, Ciência Política e Antropologia.
É nesse lugar de \textbf{fronteiras} abertas e trocas enriquecedoras \textbf{que} está a maior contribuição da área das \textbf{Humanidades} para a formação dos/as estudantes.
Elas trazem \textbf{um} diálogo não somente entre sim, mas também com as demais áreas do conhecimento, \textbf{visto} \textbf{que} \textbf{um} dos seus pressupostos \textbf{metodológico} está na motivação, no \textbf{despertar} da dúvida e da curiosidade \textbf{que} move as pesquisas científicas, na desconstrução das narrativas \textbf{que} dão suporte a visões de mundo, nas proposições e aberturas para o diálogo entre saberes de diversas sociedades, elementos essenciais para \textbf{que} a escola se torne \textbf{um} ponto de partida para a formação de \textbf{um} indivíduo consciente do seu potencial de aprendizado e motivado a atuar de forma autônoma em relação às informações e os saberes com os quais se confrontam na dinâmica da sua vida cotidiana.
Elas são assim o \textbf{que} Aristóteles chama da “ causa eficiente ”, ou seja, a causa primeira da dúvida e a curiosidade \textbf{que} \textbf{instigam} a necessidade de construir conhecimento.
Partindo da compreensão de \textbf{que} a ciência é, antes de mais nada, \textbf{uma} busca pelo saber, procuramos trazer \textbf{aqui} \textbf{uma} perspectiva de trabalho \textbf{que} sirva como auxílio aos professores e professoras das Ciências Humanas para \textbf{que} estes \textbf{instiguem} os \textbf{discentes} a trilhar pelos caminhos do saber produzido por estas ciências, fazendo uso das suas primícias de trabalho, seus objetos analíticos, suas epistemologias e com isso alcançar \textbf{uma} dinâmica de \textbf{ensino-aprendizagem} em \textbf{que} nossos/as estudantes sejam incentivados e direcionados a ter autonomia na construção de \textbf{um} pensamento crítico perante os desafios do mundo \textbf{contemporâneo}, da vida em sociedade, do seu cotidiano em todas as suas múltiplas e variadas manifestações e dimensões.
A contribuição dessas quatro ciências na construção de trajetórias de vida dos/as estudantes, nas suas experiências cotidianas, na elaboração de reflexões e ações \textbf{que} os encaminhe em \textbf{direção} a modelos sociais em \textbf{que} o processo de cidadania, de igualdade social, de respeito às diferenças culturais encontrem possibilidade de viabilizar -se, está na própria dinâmica teórico-metodológicas \textbf{que} lhes são \textbf{constitutivas}.
A relevância desse trabalho tem se mostrado cada vez mais presente na nossa sociedade, pelo teor de complexidade \textbf{que} ela \textbf{revela} e os desafios \textbf{que} as experiências sociais, políticas e culturais da \textbf{contemporaneidade} nos têm apresentado.
Nesse contexto, o professor e a professora devem promover o diálogo entre as diferentes formas de conhecimento, buscando promover encontros de saberes a partir de \textbf{uma} prática pedagógica intercultural \textbf{baseada} na horizontalidade \textbf{que} considerando o lugar de fala dos agentes sociais.
Partir de \textbf{uma} perspectiva hierárquica \textbf{que} desvaloriza as realidades e os saberes do outro impossibilitará \textbf{uma} educação \textbf{dialógica}, convertendo a prática docente em imposições de visões de mundo, o \textbf{que} é característica de \textbf{uma} \textbf{pedagogia} depositária e, portanto, opressora.
Em sociedades como a brasileira e a alagoana, tão marcadas pelas desigualdades sociais e econômicas, conhecer os mecanismos e as dinâmicas histórico-sociais \textbf{que} contribuíram e contribuem para essa realidade, desenvolvendo senso crítico autônomo \textbf{que} possibilite compreendê -la e posicionar -se frente a ela, assim como desenvolver capacidade para construir alternativas \textbf{que} escapem a esse modelo, são aportes \textbf{que} esta área do conhecimento tem a oferecer à formação dos e das estudantes.
Recentemente estas ciências têm feito reflexões sobre a necessidade de valorizar os diferentes conhecimentos, num processo \textbf{conhecido} como “ descolonização dos saberes ”.
Ainda \textbf{que} haja avanços consideráveis na busca pela compreensão ou tradução do outro, \textbf{notadamente} na Antropologia, há obstáculos interpretativos ou descritivos não superados.
A busca por compreender antologias outras ( grosso modo, as essências humanas dos outros ) é marcada por “ zonas de silêncio ”, as quais apenas esses “ outros ” as conhecem e, portanto, têm condições reais de falar delas¹⁵ ; dai a importância do “ lugar de fala ”.
Permitir \textbf{que} o outro fale possibilita \textbf{que} \textbf{que} narrativas silenciadas sejam ouvidas e compreendidas, bem como saberes outros sejam considerados no processo de ensino aprendizagem.
Nessa \textbf{direção}, práticas \textbf{que} visam promover o encontro de saberes ( como convidar representantes grupos sociais \textbf{que} geralmente não tem voz para conversar com os/as estudantes sobre temas \textbf{que} os envolvem ) devem ser consideradas no âmbito escolar.
Importa chamar para a relação entre o “ outro ” e “ nós ”.
Primeiro porque tal diferenciação não é tão clara, havendo cisões e misturas, hibridez e ambiguidade, mesclagens e atritos.
Segundo, porque o “ outro ” não resulta necessariamente de processos sociais completamente apartados de nós, em nem nós somos essencialmente “ formados ” sem eles, ainda \textbf{que} como contraposição.
Ou seja, o outro muitas das vezes não é assim tão “ outro ”.
Desta forma, compreender o outro é compreender também a sim mesmo.
O ponto fundamental \textbf{aqui} é adotar com seriedade os conhecimentos e \textbf{categorias} intelectuais de culturas de diversas sociedades, em lugar de observá -las a partir do olhar de superioridade do Ocidente, \textbf{que} as colocam como \textbf{erros} ou “ superstições ” ( BURKE : 1992 : 162 ).
Isto é de \textbf{suma} importância quando se trata da nossa experiência histórica, \textbf{visto} \textbf{que} somos \textbf{um} povo \textbf{que} se formou a partir de grupos socioculturais heterogêneos, muitos deles desconsiderados em seus saberes e experiências históricas até então.
Na prática essa mudança epistemológica significa abrir espaço para \textbf{uma} \textbf{pedagogia} intercultural, decolonial e horizontal seja praticada e os saberes “ locais ” e dos “ outros ” sejam considerados e valorizados no currículo escolar. ¹⁵ Problema \textbf{que} Eduardo Viveiro de Castro bem demonstrou em “ O nativo relativo ”, artigo publicado em 2002 na revista Mana ( CASTRO : 2002 ).

% matched lemmas: abordagem, aqui, atual, basear, categoria, conhecido, constitutivo, contemporaneidade, contemporâneo, demandar, despertar, dialético, dialógico, direção, discente, disciplina, dizer, ensino-aprendizagem, erro, forte, fronteira, hoje, humanidade, instigar, interdisciplinaridade, metodológico, método, notadamente, notar, pedagogia, que, revelar, salientar, sumo, século, teoria, territorialidade, teórico, um, uma, ver, viver
\end{textsample}
